\subsection{分离拟完备空间中的表示}

命题 7. --- 设 \(\varrho\) 是 G 在分离、局部凸且拟完备空间 E 上的连续表示。E 的有限维子表示之和在 E 中稠密。

设 \(x \in \mathrm{E}\)，U 是 \(x\) 的一个开邻域。测度 \(\nu \in \mathcal{M}(\mathrm{G})\) 满足 \(\varrho(\nu)x \in \mathrm{U}\) 的所成集合，对于紧收敛拓扑在 \(\mathcal{M}(\mathrm{G})\) 中是开的（参见 V, p. 400, no. 2）。它包含 \(\varepsilon_e\)；因此它包含形如 \(\nu = f_1 \cdot \mu\) 的测度，其中 \(f_1 \in \mathcal{C}(\mathrm{G})\)（INT, VIII, p. 171, § 4, n° 7, 命题 19）；于是它还包含形如 \(f_2 \cdot \mu\) 的测度，其中 \(f_2\) 是 G-有限函数（V, p. 464, 命题 6 及 V, p. 462, 命题 4）。

\(\gamma_{G}\) 中由 \(f_{2}\) 生成的子表示 F 是有限维的。令 \(\widetilde{F}\) 为线性映射 \(f \mapsto \varrho(f)x\) 从 F 到 E 的像。空间 \(\widetilde{F}\) 是有限维的，且包含属于 U 的元 \(\varrho(f_{2})x\)。由于 \(\varrho(g)\varrho(f) = \varrho(\gamma_{\mathrm{G}}(g)f)\) 对所有 \(g \in G\) 和每个 \(f \in F\) 成立（V, p. 406, 公式 (1)），空间 \(\widetilde{F}\) 是 \(\varrho\) 的一个与 U 相交的子表示。命题得证。

推论 1. --- 设 π 是 G 在局部凸分离拟完备空间 E 上的连续不可约表示。空间 E 是有限维的。

推论 2. --- 设 \(\varrho\) 是 G 在局部凸分离拟完备空间 E 上的连续表示，\(\pi\) 是 G 的连续不可约表示。

\begin{enumerate}
\def\labelenumi{\alph{enumi})}
\item
  从 E 到 E 的 G-态射 \(p_{\pi} = \dim(\pi)\varrho(\overline{\chi}_{\pi})\) 是 E 的连续投影，其像是 \(\pi\)-同型分量的 \(\varrho\)；
\item
  若 \(\varrho\) 是酉表示，则投影 \(p_{\pi}\) 是 E 的正交投影，其像为 \(\mathrm{M}_{\pi}(\varrho)\)。
\end{enumerate}

证明 \(a\)。线性映射 \(p_{\pi} = \dim(\pi)\varrho(\overline{\chi}_{\pi})\) 定义良好，因为 E 拟完备且 G 紧（V, p. 401）。它属于 \(\mathrm{Hom}_{\mathrm{G}}(\varrho,\varrho)\)，因为 \(\pi\) 的特征标是 G 上的连续中心函数。

设 \(\varpi\) 是 E 的有限维子表示，F 是其空间。映射 \(p_{\pi}\) 通过对子空间的传递在 F 上诱导自同态 \(\dim(\pi)\varpi(\overline{\chi}_{\pi})\)。因此，当 \(\varpi\) 不同构于 \(\pi\) 时此自同态为零，否则是 F 的恒等映射（事实上，当 \(\varpi\) 不可约时由 V, p. 460, 推论 5 得到一般情形，而一般情形由此推出，因为由 V, p. 457, 命题 1，\(\varpi\) 是半单的）。

最后，由于 E 的有限维子表示之和在 E 中稠密（命题 7），且 \(p_{\pi}\) 连续，我们得到 \(p_{\pi}\) 是一个投影，其像是 \(\pi\)-同型分量的 \(\varrho\)。若 E 是希尔伯特空间，则由 V, p. 458, 推论 1，投影 \(p_{\pi}\) 是埃尔米特的，因而是正交投影（I, p. 133, 引理 3, (ii)）。

推论 3. --- 设 \(\varrho\) 是 G 在局部凸分离拟完备空间 E 上的连续表示。自同态

\[
x \mapsto \int_ {\mathrm{G}} \varrho (g) x d \mu (g)
\]

是 E 的一个投影，其像是 E 中不变向量的空间 \(E^{G}\)。特别地，若 E 有限维，则

\[
\mathrm{dim} (\mathrm{E} ^ {\mathrm{G}}) = \int_ {\mathrm{G}} \chi_ {\varrho} (g) d \mu (g).
\]

第一个断言是前一推论在 \(\pi\) 为 G 的一维平凡表示时的特殊情形。第二个断言由此得出，因为 \(E^{G}\) 的维数是到 \(E^{G}\) 上的正交投影的迹，即

\[
\mathrm{dim} (\mathrm{E} ^ {\mathrm{G}}) = \mathrm{Tr} \Bigl (\int_ {\mathrm{G}} \varrho (g) d \mu (g) \Bigr) = \int_ {\mathrm{G}} \chi_ {\varrho} (g) d \mu (g).
\]

特别地，当 E 有限维时，E 中存在向量 \(x \neq 0\)，使得 \(\varrho(g)x = x\) 对所有 \(g \in G\) 成立，当且仅当

\[
\int_ {\mathrm{G}} \chi_ {\varrho} (g) d \mu (g) \neq 0.
\]

推论 4. --- 设 \(\varrho\) 是 G 在希尔伯特空间 E 上的酉表示。空间 E 是各 \(\pi\)-同型分量 \(\mathrm{M}_{\pi}(\varrho)\)（其中 \(\pi \in \widehat{\mathrm{G}}\)）的希尔伯特直和。特别地，表示 \(\varrho\) 是半单的。

对应于 G 的不同构不可约表示的 \(\varrho\) 的同型分量彼此正交（V, p. 394, 命题 8）。由于 G 的每个有限维酉表示都是半单的（V, p. 457, 命题 1），各同型分量 \(M_{\pi}(\varrho)\) 对 \(\pi \in \widehat{G}\) 的和在 E 中稠密（命题 7）。推论由此得出。
% semantic-fragments: legacy-input-seam
\relax{}
